problem:  
Let \((\mathfrak n,\langle\cdot,\cdot\rangle)\) be a fixed **nilsoliton**—a nilpotent Lie algebra equipped with its canonical (up to scaling and automorphism) nilsoliton metric satisfying  
\[
\mathrm{Ric} = c\, I + D_E^s, \qquad c < 0,
\]
where \(D_E^s\) is the symmetric **Einstein derivation**.  
Consider all solvable Lie algebras of the form  
\[
\mathfrak s = \mathfrak a \ltimes \mathfrak n, \qquad \text{with }\ [\mathfrak a, \mathfrak a] = 0,
\]
such that each \(A \in \mathfrak a\) acts on \(\mathfrak n\) by a symmetric derivation commuting with \(D_E^s\). Each such choice defines a solvsoliton metric via Lauret’s construction.  

**Question:**  
For a fixed nilsoliton pair \((\mathfrak n, D_E^s)\), classify the possible **commuting families** of symmetric derivations on \(\mathfrak n\) that yield distinct (up to scaling and isometry) solvsoliton extensions.  
Specifically, determine whether the space of such commuting families (modulo the automorphism group \(\mathrm{Aut}(\mathfrak n)\)) is discrete, finite‑dimensional, or admits continuous moduli.  
If it admits continuous moduli, characterize the deformation space: does varying the abelian derivation family change only the extension algebra, or can it produce inequivalent solvsoliton metrics on diffeomorphic solvmanifolds?

Why is it a "good" problem:  
This formulation sits precisely at the known frontier of solvsoliton theory. It builds on the established uniqueness of nilsolitons and isolates the true open phenomenon—the spectrum and geometry of **commuting symmetric derivations** compatible with a fixed nilsoliton structure. Understanding whether these extensions form a discrete or genuinely continuous moduli space probes the rigidity of homogeneous Ricci solitons at their algebraic core. A resolution would link algebraic deformation theory of Lie algebras with Ricci flow dynamics, revealing deeper structural constraints or flexibility in homogeneous Einstein-like geometries. The statement is concise, well‑posed, and demands interplay between Lie algebra cohomology, derivation algebras, and geometric analysis, embodying the qualities of a "good" research-level problem.